\section{Motivação}

A medida que a ciência avança e novas teorias físicas vão sendo desenvolvidas, é necessário refinar cada vez mais a técnicas experimentais para comprovar ou refutar tais teorias.
Uma dessas técnicas é a interferometria, que consiste na utilização do fenômeno da interferência da luz para realizar medidas.
Com ela, é possível medir distâncias da ordem de grandeza do comprimento de onda da luz utilizada. 
Um exemplo disso foi quando, em 1887, A. Michelson e E. Morley realizaram seu famoso experimento para determinar a existência do éter luminífero.
O aparato que os dois cientistas usaram foi um interferômetro a laser desenvolvido por Michelson, que recebeu seu nome. 
Assim, eles demonstraram a inexistência do éter, e o interferômetro de Michelson foi utilizado por diversos experimentos ao longo dos últimos dois śeculos, tendo sido responsável pela detecção de ondas gravitácionais em 2014 \ccite{LIGO}.

Neste trabalho, buscamos constatar a eficácia e a precisão do interferômetro de Michelson ao medir 3 grandezas ópticas. 
Na parte 1, buscamos determinar o comprimento de onda do laser de HeNe; 
na parte 2, buscamos determinar como o índice de refração do ar varia com a pressão; 
e na parte 3, buscamos determinar o índice de refração de um vidro.
